\documentclass[a4paper,11pt]{report}

\usepackage[utf8]{inputenc}
\usepackage[T1]{fontenc}
\usepackage[spanish,es-tabla,es-noquoting]{babel}
\usepackage{geometry}
\usepackage{array}
\usepackage{booktabs}
\usepackage{listings}
\usepackage{tikz}
\usepackage{hyperref}
\usetikzlibrary{positioning,arrows.meta}

\geometry{margin=2.5cm}
\setlength{\emergencystretch}{2em}
\renewcommand{\chaptername}{Capítulo}
\renewcommand{\contentsname}{Índice}
\renewcommand{\listtablename}{Índice de tablas}
\renewcommand{\lstlistingname}{Listado}

\lstset{
    basicstyle=\small\ttfamily,
    columns=fullflexible,
    keepspaces=true,
    breaklines=true,
    frame=single,
    numbers=left,
    numberstyle=\tiny,
    showstringspaces=false,
    tabsize=4,
    literate=
        {á}{{\'a}}1 {é}{{\'e}}1 {í}{{\'{\i}}}1 {ó}{{\'o}}1 {ú}{{\'u}}1
        {ñ}{{\~n}}1
}

\title{Práctica 2\\Listas doblemente enlazadas circulares}
\author{Samuel Carrasco Santurde\\Alexander Tomás Echevarria Pizarro\\Erasmo Seebold Guadilla}
\date{}

\begin{document}
\chapter{Aspectos pendientes para completar y mejorar el proyecto}


\section{Corregir las operaciones con casos límite}
\begin{enumerate}
  \item \textcolor{blue}{Corregido; falta prueba específica} \textbf{Evitar el error de \texttt{contains}.} La implementación comprueba si \texttt{locate(elem)} devuelve \texttt{null} y devuelve \texttt{false} si el elemento no existe. Añadir pruebas para lista vacía y elemento ausente.
  \item \textcolor{blue}{Corregido y probado parcialmente} \textbf{Definir qué ocurre si falta el objetivo de \texttt{addAfter}.} Si la lista no está vacía y \texttt{target} no está presente, no se inserta nada. Hay pruebas para objetivo presente y ausente; falta probar la lista vacía.
  \item \textbf{Pendiente: proteger y probar la auto-fusión.} El método \texttt{merge} no comprueba explícitamente si la lista recibida es la misma lista destino. Definir el comportamiento y añadir una prueba.
\end{enumerate}
\section{Completar y validar los merges ordenados}

\begin{enumerate}
    \item \textcolor{blue}{Implementado; falta probar} \textbf{Resolver la igualdad en \texttt{merge}.} El método combina elementos iguales y avanza ambos punteros. Faltan pruebas que verifiquen que no hay duplicados y que las relaciones asociadas se combinan correctamente.
    \item \textcolor{blue}{Retirado} \textbf{Eliminar \texttt{mergePruebas}.} La variante experimental ya no está en \texttt{OrderedDouble\allowbreak LinkedList}; no requiere pruebas.
    \item \textbf{Pendiente: documentar el comportamiento de la lista fuente.} \texttt{merge} conserva la fuente cuando el destino no está vacío, pero comparte sus nodos si el destino está vacío. Definir si ese caso debe copiar o transferir nodos y añadir pruebas.
    \item \textbf{Pendiente: probar \texttt{merge}.} Cubrir listas vacías, rangos intercalados y separados, elementos iguales en distintas posiciones, tamaños diferentes y auto-fusión; comprobar orden, tamaño, duplicados y estado de la fuente.
    \item \textbf{Pendiente: verificar las relaciones de actores y películas.} Comprobar que los IDs no se duplican y que cada relación existe en ambas direcciones, también al combinar lotes.
\end{enumerate}

\section{Aumentar la cobertura de pruebas}

Hay dos pruebas de \texttt{addAfter} y una prueba de lectura de ficheros. Esta última comprueba cantidades de registros, actores y películas, algunos actores y la búsqueda de un actor por ID. Aún faltan pruebas automatizadas para listas vacías, de uno y varios elementos; extremos e interiores; elementos presentes y ausentes; inserciones repetidas; y eliminaciones. En listas ordenadas, comprobar también orden, tamaño y duplicados según \texttt{compareTo}.

Para los merges conservados, probar listas fuente y destino vacías, elementos intercalados, rangos sin solapamiento, elementos iguales y listas de tamaños distintos. Verificar si la fuente se conserva y probar la auto-fusión.

En el ejemplo de actores y películas, comprobar que las cantidades de entidades son las esperadas, que no se repite ningún ID y que cada relación que aparece en la lista de películas de un actor aparece también en la lista de actores de esa película, y viceversa.

\section{Alinear el diseño con el enunciado}

\begin{enumerate}
    \item \textbf{Revisar el nodo centinela.} El enunciado describe una lista circular con centinela, pero nuestra implementación representa la lista vacía con \texttt{first == null}. Decidir con el profesorado si se debe cambiar la implementación para usar centinela o mantener esta variante y explicar claramente la diferencia.
    \item \textbf{Revisar el carácter de \texttt{DoubleLinkedList}.} El enunciado la presenta como clase abstracta, mientras que en el proyecto es concreta. Si la abstracción es un requisito de evaluación, modificar la declaración y confirmar que las clases derivadas cubren el contrato.
    \item \textbf{Separar, si se permite, el dominio de la lista genérica.} \texttt{OrderedDoubleLinkedList} conoce actualmente a \texttt{Actor} y \texttt{Pelicula} para combinar elementos iguales. Una mejora sería dejar que la lista gestione solo orden e inserción y que la lógica de combinación pertenezca a las clases de dominio o al lector, siempre respetando las herramientas permitidas en la práctica.
\end{enumerate}

\section{Mejorar las mediciones de eficiencia}

El lector de listas enlazadas usa actualmente \texttt{merge} y procesa los registros en lotes de 127; los parámetros para seleccionar métodos se han retirado. No se está comparando más de una estrategia desde este lector. Si se mantienen mediciones de eficiencia, ejecutarlas con datos equivalentes y listas nuevas, repetirlas y distinguir lectura de fusión.

El cronómetro informa milisegundos, pero usa \texttt{System.currentTimeMillis()}; para intervalos breves convendría un reloj de mayor resolución y varias ejecuciones antes de extraer conclusiones.

\section{Criterio para dar el trabajo por terminado}

Consideraríamos el proyecto listo cuando las operaciones respeten sus contratos e invariantes, el método de fusión conservado no pierda ni duplique elementos, las relaciones del ejemplo estén sincronizadas, las pruebas cubran los casos normales y límite, y las mediciones que se presenten sean reproducibles. Cualquier diferencia intencionada respecto al enunciado debe quedar explicada en la memoria.

\textcolor{blue}{Actualización: se ha eliminado \texttt{mergePruebas} y el lector de listas enlazadas ahora utiliza \texttt{merge}.}
\end{document}
